\documentclass[12pt,a4paper]{article}

\usepackage[T2A]{fontenc}
\usepackage[utf8]{inputenc}
\usepackage[russian]{babel}
\usepackage{amsmath,amssymb,amsthm,mathtools}
\usepackage{enumitem}
\usepackage{array}
\usepackage{booktabs}
\usepackage{hyperref}
\usepackage{listings}
\usepackage{xcolor}

\hypersetup{
    colorlinks=true,
    linkcolor=blue,
    urlcolor=blue
}

\newtheorem{definition}{Определение}
\newtheorem{theorem}{Теорема}
\newtheorem{lemma}{Лемма}
\newtheorem{example}{Пример}

\DeclareMathOperator{\Enc}{Enc}
\DeclareMathOperator{\Dec}{Dec}
\DeclareMathOperator{\Gen}{Gen}
\DeclareMathOperator{\Sign}{Sign}
\DeclareMathOperator{\Ver}{Ver}
\DeclareMathOperator{\Mac}{Mac}
\DeclareMathOperator{\KeyGen}{KeyGen}
\DeclareMathOperator{\Hash}{Hash}
\DeclareMathOperator{\Adv}{Adv}

\title{Билет №16 \\ \small Основы криптографии. Задачи обеспечения конфиденциальности и целостности информации. Системы шифрования с открытым ключом (RSA). Цифровая подпись. (ИТМО) \\ Основы криптографии. Задачи обеспечения конфиденциальности и целостности информации. Теоретико-информационный и теоретико-сложностный подходы к определению криптографической стойкости. Стандарты шифрования DES и AES, современные российские алгоритмы шифрования ГОСТ. Системы шифрования с открытым ключом (RSA). Цифровая подпись. Методы генерации и распределения ключей. (СпБГУ)}
\author{Сериков Владислав Александрович}
\date{5 мая 2026}

\begin{document}

\maketitle
\tableofcontents
\newpage

\section{Предмет криптографии и базовые задачи защиты информации}

\subsection{Криптография, криптоанализ, криптология}

\begin{definition}
\textbf{Криптография} --- область математики и информатики, изучающая методы защиты информации с помощью преобразований, зависящих от секретных параметров, называемых ключами.
\end{definition}

\begin{definition}
\textbf{Криптоанализ} --- область, изучающая методы нарушения криптографической защиты: восстановление открытого текста, ключей, подделку сообщений, поиск коллизий и т.д.
\end{definition}

\begin{definition}
\textbf{Криптология} --- совокупность криптографии и криптоанализа.
\end{definition}

Современная криптография решает не только задачу сокрытия сообщений, но и задачи проверки целостности, аутентификации, невозможности отказа от авторства, безопасного обмена ключами, построения защищённых протоколов.

\subsection{Основные цели защиты информации}

Основные свойства защищённой информации:

\begin{enumerate}
    \item \textbf{Конфиденциальность} --- информация доступна только тем субъектам, которым она предназначена.
    \item \textbf{Целостность} --- информация не была изменена несанкционированным образом.
    \item \textbf{Аутентичность} --- возможность проверить источник сообщения или личность участника протокола.
    \item \textbf{Неотказуемость} --- отправитель не может обоснованно отрицать факт создания или подписания сообщения.
    \item \textbf{Доступность} --- система остаётся пригодной для использования авторизованными пользователями.
\end{enumerate}

В данном конспекте основной акцент делается на конфиденциальности и целостности.

\section{Математическая модель шифрования}

\subsection{Симметрическая криптосистема}

\begin{definition}
\textbf{Симметрическая криптосистема} задаётся тройкой вероятностных или детерминированных алгоритмов
\[
    \Pi = (\Gen, \Enc, \Dec),
\]
где:
\begin{itemize}
    \item $\Gen(1^n)$ --- алгоритм генерации секретного ключа $k$;
    \item $\Enc_k(m)$ --- алгоритм шифрования сообщения $m$ на ключе $k$;
    \item $\Dec_k(c)$ --- алгоритм расшифрования шифртекста $c$ на ключе $k$.
\end{itemize}
\end{definition}

Корректность требует, чтобы для любого допустимого сообщения $m$ и ключа $k$, порождённого $\Gen$,
\[
    \Dec_k(\Enc_k(m)) = m.
\]

В симметрической криптографии один и тот же секретный ключ используется для шифрования и расшифрования.

\subsection{Асимметрическая криптосистема}

\begin{definition}
\textbf{Криптосистема с открытым ключом} задаётся тройкой алгоритмов
\[
    \Pi = (\KeyGen, \Enc, \Dec),
\]
где:
\begin{itemize}
    \item $\KeyGen(1^n)$ порождает пару ключей $(pk,sk)$;
    \item $pk$ --- открытый ключ, доступный всем;
    \item $sk$ --- закрытый ключ, известный только владельцу;
    \item $\Enc_{pk}(m)$ --- шифрование на открытом ключе;
    \item $\Dec_{sk}(c)$ --- расшифрование на закрытом ключе.
\end{itemize}
\end{definition}

Корректность:
\[
    \Dec_{sk}(\Enc_{pk}(m)) = m.
\]

\section{Конфиденциальность}

\subsection{Идея конфиденциальности}

Конфиденциальность означает, что противник, перехвативший шифртекст, не должен получить существенной информации об открытом тексте.

Существуют разные уровни формализации конфиденциальности:

\begin{enumerate}
    \item \textbf{Теоретико-информационная стойкость}: противник не получает информации даже при неограниченных вычислительных ресурсах.
    \item \textbf{Теоретико-сложностная стойкость}: противник ограничен полиномиальным временем, и нарушение защиты требует решения вычислительно трудной задачи.
\end{enumerate}

\section{Целостность и аутентификация сообщений}

\subsection{Коды аутентификации сообщений}

\begin{definition}
\textbf{Код аутентификации сообщения} или \textbf{MAC} --- это схема
\[
    (\Gen,\Mac,\Ver),
\]
где:
\begin{itemize}
    \item $\Gen$ генерирует секретный ключ $k$;
    \item $\Mac_k(m)$ вычисляет тег $t$ для сообщения $m$;
    \item $\Ver_k(m,t)$ возвращает $1$, если тег корректен, и $0$ иначе.
\end{itemize}
\end{definition}

Корректность:
\[
    \Ver_k(m,\Mac_k(m)) = 1.
\]

MAC обеспечивает целостность и аутентичность в симметрической модели, но не обеспечивает неотказуемость, поскольку обе стороны знают один и тот же ключ.

\subsection{Криптографические хеш-функции}

\begin{definition}
\textbf{Криптографическая хеш-функция} --- функция
\[
    H : \{0,1\}^* \to \{0,1\}^n,
\]
отображающая сообщения произвольной длины в строки фиксированной длины.
\end{definition}

Основные свойства:

\begin{enumerate}
    \item \textbf{Стойкость к нахождению прообраза}: по $h$ трудно найти $m$ такое, что $H(m)=h$.
    \item \textbf{Стойкость ко второму прообразу}: по $m$ трудно найти $m' \ne m$ такое, что $H(m')=H(m)$.
    \item \textbf{Стойкость к коллизиям}: трудно найти пару $m \ne m'$ такую, что
    \[
        H(m)=H(m').
    \]
\end{enumerate}

Хеш-функции используются в цифровых подписях, MAC, хранении паролей, построении протоколов.

\section{Теоретико-информационный подход к криптографической стойкости}

\subsection{Совершенная секретность}

Пусть $M$ --- случайная величина, описывающая открытый текст, $C$ --- шифртекст, $K$ --- ключ.

\begin{definition}
Криптосистема обладает \textbf{совершенной секретностью}, если для любых сообщений $m$ и шифртекстов $c$, для которых $\Pr[C=c]>0$, выполняется
\[
    \Pr[M=m \mid C=c] = \Pr[M=m].
\]
\end{definition}

Иными словами, наблюдение шифртекста не меняет априорного распределения сообщений.

Эквивалентная формулировка:
\[
    \Pr[C=c \mid M=m_0] = \Pr[C=c \mid M=m_1]
\]
для любых допустимых сообщений $m_0,m_1$ и любого $c$.

\subsection{Одноразовый блокнот}

\begin{definition}
\textbf{Одноразовый блокнот} --- симметрическая схема шифрования, где сообщение $m \in \{0,1\}^n$ шифруется ключом $k \in \{0,1\}^n$, выбранным равновероятно:
\[
    c = m \oplus k.
\]
Расшифрование:
\[
    m = c \oplus k.
\]
\end{definition}

\subsubsection*{Алгоритм одноразового блокнота}

\textbf{Генерация ключа:}
\begin{enumerate}
    \item Выбрать случайную строку $k \in \{0,1\}^n$.
    \item Передать ключ отправителю и получателю по защищённому каналу.
\end{enumerate}

\textbf{Шифрование:}
\begin{enumerate}
    \item Взять сообщение $m \in \{0,1\}^n$.
    \item Вычислить $c=m\oplus k$.
    \item Отправить $c$.
\end{enumerate}

\textbf{Расшифрование:}
\begin{enumerate}
    \item Получить $c$.
    \item Вычислить $m=c\oplus k$.
\end{enumerate}

\begin{example}
Пусть
\[
    m=1011,\qquad k=0110.
\]
Тогда
\[
    c=m\oplus k=1011\oplus 0110=1101.
\]
Расшифрование:
\[
    c\oplus k=1101\oplus 0110=1011=m.
\]
\end{example}

\begin{theorem}[Совершенная секретность одноразового блокнота]
Одноразовый блокнот обладает совершенной секретностью, если ключ выбирается равновероятно из $\{0,1\}^n$, используется только один раз и имеет ту же длину, что и сообщение.
\end{theorem}

\begin{proof}
Пусть $m,c \in \{0,1\}^n$ --- произвольные сообщение и шифртекст. Для того чтобы при шифровании сообщения $m$ получился шифртекст $c$, ключ должен быть равен
\[
    k = m \oplus c.
\]
Такой ключ существует и единственен. Поскольку ключ выбирается равновероятно из множества $\{0,1\}^n$, имеющего $2^n$ элементов, получаем
\[
    \Pr[C=c \mid M=m] = \Pr[K=m\oplus c] = 2^{-n}.
\]
Это значение не зависит от $m$. Следовательно, для любых $m_0,m_1$:
\[
    \Pr[C=c \mid M=m_0] = \Pr[C=c \mid M=m_1].
\]
Значит, шифртекст не даёт информации о сообщении. Поэтому схема обладает совершенной секретностью.
\end{proof}

\subsection{Теорема Шеннона о длине ключа}

\begin{theorem}[Необходимое условие совершенной секретности]
Пусть криптосистема обладает совершенной секретностью, множество сообщений $M$, множество ключей $K$ и множество шифртекстов $C$ конечны, причём каждое сообщение имеет положительную вероятность. Тогда
\[
    |K| \ge |M|.
\]
\end{theorem}

\begin{proof}
Зафиксируем некоторый шифртекст $c$, который может быть получен с ненулевой вероятностью.

Для каждого сообщения $m \in M$ из совершенной секретности следует, что
\[
    \Pr[M=m \mid C=c] = \Pr[M=m] > 0.
\]
Следовательно,
\[
    \Pr[M=m \land C=c] > 0.
\]
Значит, существует хотя бы один ключ $k \in K$, такой что
\[
    \Enc_k(m)=c.
\]

Рассмотрим отображение, которое каждому сообщению $m$ сопоставляет некоторый ключ $k_m$, переводящий $m$ в фиксированный шифртекст $c$.

Если два различных сообщения $m_1 \ne m_2$ переводятся одним и тем же ключом $k$ в один и тот же шифртекст $c$, то при расшифровании возникнет противоречие:
\[
    \Dec_k(c)=m_1
\]
и одновременно
\[
    \Dec_k(c)=m_2.
\]
Но алгоритм расшифрования должен быть корректным и однозначным. Поэтому разным сообщениям соответствуют разные ключи.

Следовательно, ключей должно быть не меньше, чем сообщений:
\[
    |K| \ge |M|.
\]
\end{proof}

Из этой теоремы следует важный вывод: совершенная секретность требует ключа не короче сообщения. Поэтому одноразовый блокнот практически неудобен для больших объёмов данных.

\section{Теоретико-сложностный подход к стойкости}

\subsection{Идея вычислительной стойкости}

В современной криптографии обычно предполагается, что противник ограничен вычислительными ресурсами. Схема считается стойкой, если любой вероятностный полиномиальный алгоритм имеет лишь пренебрежимо малую вероятность успеха.

\begin{definition}
Функция $\mu(n)$ называется \textbf{пренебрежимо малой}, если для любого полинома $p(n)$ существует $N$, такое что для всех $n>N$
\[
    \mu(n) < \frac{1}{p(n)}.
\]
\end{definition}

Примеры пренебрежимо малых функций:
\[
    2^{-n},\qquad 2^{-\sqrt n}.
\]

Не являются пренебрежимо малыми:
\[
    \frac{1}{n},\qquad \frac{1}{n^{100}}.
\]

\subsection{Безопасность через игры}

Криптографические определения часто задаются через эксперимент между противником и ``оракулом''.

Например, для шифрования используется идея неразличимости: противник выбирает два сообщения $m_0,m_1$, получает шифрование одного из них и должен угадать, какое именно было зашифровано.

\subsection{IND-CPA}

\begin{definition}
Схема шифрования обладает \textbf{стойкостью к атаке с выбранным открытым текстом} или \textbf{IND-CPA}, если никакой вероятностный полиномиальный противник не может отличить шифрование $m_0$ от шифрования $m_1$ с преимуществом, существенно отличным от нуля.
\end{definition}

Эксперимент:

\begin{enumerate}
    \item Генерируется ключ $k$.
    \item Противник выбирает два сообщения $m_0,m_1$ одинаковой длины.
    \item Случайно выбирается бит $b \in \{0,1\}$.
    \item Противнику передаётся
    \[
        c = \Enc_k(m_b).
    \]
    \item Противник выдаёт бит $b'$.
    \item Он выигрывает, если $b'=b$.
\end{enumerate}

Преимущество противника:
\[
    \Adv = \left|\Pr[b'=b]-\frac12\right|.
\]

Схема IND-CPA-стойкая, если $\Adv$ пренебрежимо мало.

\section{Симметрическое шифрование}

\subsection{Блочные и поточные шифры}

\begin{definition}
\textbf{Блочный шифр} --- семейство перестановок
\[
    E_k : \{0,1\}^n \to \{0,1\}^n,
\]
параметризованное ключом $k$.
\end{definition}

Блочный шифр шифрует блоки фиксированной длины. Для сообщений произвольной длины используются режимы работы: ECB, CBC, CFB, OFB, CTR, GCM и другие.

\begin{definition}
\textbf{Поточный шифр} генерирует псевдослучайную гамму $z_1,z_2,\ldots$, которая складывается с открытым текстом:
\[
    c_i = m_i \oplus z_i.
\]
\end{definition}

\subsection{Сеть Фейстеля}

Многие блочные шифры, включая DES и ``Магму'', основаны на сети Фейстеля.

Пусть блок делится на две половины:
\[
    X = (L_0,R_0).
\]

Один раунд сети Фейстеля:
\[
    L_i = R_{i-1},
\]
\[
    R_i = L_{i-1} \oplus F(R_{i-1},K_i).
\]

Здесь $F$ --- раундовая функция, $K_i$ --- раундовый ключ.

\subsubsection*{Обратимость сети Фейстеля}

Даже если функция $F$ необратима, весь раунд Фейстеля обратим.

Из формул
\[
    L_i=R_{i-1},\qquad R_i=L_{i-1}\oplus F(R_{i-1},K_i)
\]
получаем:
\[
    R_{i-1}=L_i,
\]
\[
    L_{i-1}=R_i\oplus F(L_i,K_i).
\]

\begin{theorem}
Раунд сети Фейстеля является обратимым преобразованием для любой функции $F$.
\end{theorem}

\begin{proof}
Пусть после раунда получены $(L_i,R_i)$. Тогда исходные значения восстанавливаются по формулам:
\[
    R_{i-1}=L_i,
\]
\[
    L_{i-1}=R_i\oplus F(L_i,K_i).
\]
Поскольку операция XOR обратима сама себе, восстановление однозначно. Следовательно, преобразование обратимо.
\end{proof}

\section{DES}

\subsection{Общие сведения}

DES, Data Encryption Standard, --- блочный симметрический шифр. Основные параметры:

\begin{center}
\begin{tabular}{ll}
\toprule
Размер блока & 64 бита \\
Эффективная длина ключа & 56 бит \\
Структура & 16-раундовая сеть Фейстеля \\
Размер половины блока & 32 бита \\
\bottomrule
\end{tabular}
\end{center}

Формально ключ DES имеет 64 бита, но 8 бит используются как биты контроля чётности, поэтому эффективная длина ключа равна 56 битам.

На современном уровне DES считается устаревшим из-за малой длины ключа. Его развитие --- Triple DES, но и он постепенно выводится из практического применения.

\subsection{Общая схема DES}

Алгоритм DES:

\begin{enumerate}
    \item Применить начальную перестановку $IP$ к 64-битному блоку.
    \item Разделить результат на две 32-битные половины:
    \[
        L_0,R_0.
    \]
    \item Выполнить 16 раундов Фейстеля:
    \[
        L_i = R_{i-1},
    \]
    \[
        R_i = L_{i-1} \oplus f(R_{i-1},K_i).
    \]
    \item Поменять половины местами.
    \item Применить конечную перестановку $IP^{-1}$.
\end{enumerate}

\subsection{Раундовая функция DES}

Функция $f$ устроена так:

\begin{enumerate}
    \item \textbf{Расширение}: 32-битное значение $R$ расширяется до 48 бит с помощью перестановки-расширения $E$.
    \[
        E(R) \in \{0,1\}^{48}.
    \]
    \item \textbf{Сложение с раундовым ключом}:
    \[
        B = E(R)\oplus K_i.
    \]
    \item \textbf{S-блоки}: 48 бит разбиваются на восемь 6-битных блоков. Каждый S-блок преобразует 6 бит в 4 бита:
    \[
        S_j:\{0,1\}^6\to \{0,1\}^4.
    \]
    \item \textbf{Перестановка}: полученные 32 бита переставляются по фиксированной таблице $P$.
\end{enumerate}

\subsection{Минимальная иллюстрация раунда Фейстеля}

Пусть для учебного примера блок имеет 8 бит:
\[
    L_0=1010,\qquad R_0=0110.
\]
Пусть
\[
    F(R_0,K_1)=1100.
\]
Тогда:
\[
    L_1=R_0=0110,
\]
\[
    R_1=L_0\oplus F(R_0,K_1)=1010\oplus 1100=0110.
\]
Итак,
\[
    (L_1,R_1)=(0110,0110).
\]

\section{AES}

\subsection{Общие сведения}

AES, Advanced Encryption Standard, --- современный стандарт блочного симметрического шифрования.

Основные параметры:

\begin{center}
\begin{tabular}{llll}
\toprule
Вариант & Размер блока & Длина ключа & Число раундов \\
\midrule
AES-128 & 128 бит & 128 бит & 10 \\
AES-192 & 128 бит & 192 бит & 12 \\
AES-256 & 128 бит & 256 бит & 14 \\
\bottomrule
\end{tabular}
\end{center}

AES не является сетью Фейстеля. Он основан на substitution-permutation network, то есть сети замен и перестановок.

\subsection{Представление состояния}

128-битный блок представляется как матрица $4\times 4$ байта:
\[
    S =
    \begin{pmatrix}
    s_{0,0} & s_{0,1} & s_{0,2} & s_{0,3}\\
    s_{1,0} & s_{1,1} & s_{1,2} & s_{1,3}\\
    s_{2,0} & s_{2,1} & s_{2,2} & s_{2,3}\\
    s_{3,0} & s_{3,1} & s_{3,2} & s_{3,3}
    \end{pmatrix}.
\]

Операции AES выполняются над байтами, которые рассматриваются как элементы поля
\[
    GF(2^8).
\]

\subsection{Раунды AES}

AES состоит из следующих преобразований:

\begin{enumerate}
    \item \textbf{SubBytes} --- нелинейная байтовая замена с помощью S-блока.
    \item \textbf{ShiftRows} --- циклический сдвиг строк матрицы состояния.
    \item \textbf{MixColumns} --- линейное перемешивание байтов внутри каждого столбца.
    \item \textbf{AddRoundKey} --- XOR состояния с раундовым ключом.
\end{enumerate}

\subsection{Алгоритм AES-128}

\begin{enumerate}
    \item Сгенерировать 11 раундовых ключей из исходного 128-битного ключа.
    \item Выполнить начальный шаг:
    \[
        S \leftarrow S \oplus K_0.
    \]
    \item Для раундов $1,\ldots,9$ выполнить:
    \begin{enumerate}
        \item SubBytes.
        \item ShiftRows.
        \item MixColumns.
        \item AddRoundKey.
    \end{enumerate}
    \item В последнем, десятом, раунде выполнить:
    \begin{enumerate}
        \item SubBytes.
        \item ShiftRows.
        \item AddRoundKey.
    \end{enumerate}
\end{enumerate}

\subsection{Иллюстрация ShiftRows}

Пусть состояние имеет вид:
\[
\begin{pmatrix}
a_0&a_1&a_2&a_3\\
b_0&b_1&b_2&b_3\\
c_0&c_1&c_2&c_3\\
d_0&d_1&d_2&d_3
\end{pmatrix}.
\]

После ShiftRows:
\[
\begin{pmatrix}
a_0&a_1&a_2&a_3\\
b_1&b_2&b_3&b_0\\
c_2&c_3&c_0&c_1\\
d_3&d_0&d_1&d_2
\end{pmatrix}.
\]

\subsection{Сравнение DES и AES}

\begin{center}
\begin{tabular}{lll}
\toprule
Параметр & DES & AES \\
\midrule
Размер блока & 64 бита & 128 бит \\
Длина ключа & 56 бит & 128/192/256 бит \\
Структура & сеть Фейстеля & SP-сеть \\
Число раундов & 16 & 10/12/14 \\
Практический статус & устарел & широко применяется \\
\bottomrule
\end{tabular}
\end{center}

\section{Современные российские алгоритмы шифрования ГОСТ}

\subsection{ГОСТ Р 34.12-2015 и ГОСТ 34.12-2018}

Современный российский стандарт блочного шифрования описывает два алгоритма:

\begin{enumerate}
    \item \textbf{Магма} --- блочный шифр с размером блока 64 бита и ключом 256 бит.
    \item \textbf{Кузнечик} --- блочный шифр с размером блока 128 бит и ключом 256 бит.
\end{enumerate}

\subsection{Магма}

\subsubsection{Общие параметры}

\begin{center}
\begin{tabular}{ll}
\toprule
Размер блока & 64 бита \\
Размер ключа & 256 бит \\
Структура & сеть Фейстеля \\
Число раундов & 32 \\
Размер половины блока & 32 бита \\
\bottomrule
\end{tabular}
\end{center}

\subsubsection{Раунд Магмы}

Пусть блок разделён на две 32-битные половины:
\[
    (a_1,a_0).
\]

Один раунд имеет вид:
\[
    a_1' = a_0,
\]
\[
    a_0' = a_1 \oplus g(a_0,k_i).
\]

Функция $g$ включает:

\begin{enumerate}
    \item Сложение по модулю $2^{32}$:
    \[
        x = (a_0+k_i)\bmod 2^{32}.
    \]
    \item Нелинейную замену с помощью S-блоков.
    \item Циклический сдвиг влево на 11 бит.
\end{enumerate}

\subsubsection{Упрощённый учебный пример}

Пусть половины имеют по 4 бита:
\[
    a_1=1010,\qquad a_0=0011.
\]
Пусть учебная функция
\[
    g(a_0,k_i)=1100.
\]
Тогда:
\[
    a_1'=0011,
\]
\[
    a_0'=1010\oplus 1100=0110.
\]
Новый блок:
\[
    (a_1',a_0')=(0011,0110).
\]

\subsection{Кузнечик}

\subsubsection{Общие параметры}

\begin{center}
\begin{tabular}{ll}
\toprule
Размер блока & 128 бит \\
Размер ключа & 256 бит \\
Структура & substitution-linear network \\
Число раундов & 10 \\
\bottomrule
\end{tabular}
\end{center}

\subsubsection{Основные преобразования}

Кузнечик использует три основных преобразования:

\begin{enumerate}
    \item $X[k]$ --- сложение с ключом:
    \[
        X[k](a)=a\oplus k.
    \]
    \item $S$ --- нелинейная байтовая замена.
    \item $L$ --- линейное преобразование.
\end{enumerate}

Раундовое преобразование можно записать как:
\[
    LSX[k_i](a)=L(S(a\oplus k_i)).
\]

Шифрование:
\[
    E_{K}(a)=X[k_{10}]\,LSX[k_9]\cdots LSX[k_1](a).
\]

То есть первые 9 раундов включают XOR с ключом, замену и линейное преобразование, а в конце выполняется добавление последнего ключа.

\section{Асимметрическая криптография}

\subsection{Идея открытого ключа}

В симметрической криптографии отправитель и получатель должны заранее иметь общий секретный ключ. Это создаёт проблему распределения ключей.

В криптографии с открытым ключом у каждого участника есть:

\begin{itemize}
    \item открытый ключ $pk$, который можно публиковать;
    \item закрытый ключ $sk$, который должен храниться в секрете.
\end{itemize}

Открытый ключ используется для шифрования или проверки подписи, закрытый --- для расшифрования или создания подписи.

\section{RSA}

\subsection{Математические основы RSA}

RSA основан на арифметике по модулю
\[
    n=pq,
\]
где $p$ и $q$ --- большие простые числа.

Функция Эйлера:
\[
    \varphi(n)=\varphi(pq)=(p-1)(q-1).
\]

\begin{theorem}[Теорема Эйлера]
Если $\gcd(a,n)=1$, то
\[
    a^{\varphi(n)}\equiv 1 \pmod n.
\]
\end{theorem}

\begin{proof}
Рассмотрим множество обратимых классов по модулю $n$:
\[
    U_n=\{x \in \mathbb{Z}_n : \gcd(x,n)=1\}.
\]
Количество элементов этого множества равно $\varphi(n)$.

Так как $\gcd(a,n)=1$, умножение на $a$ по модулю $n$ задаёт перестановку множества $U_n$:
\[
    x \mapsto ax \bmod n.
\]
Следовательно, произведение всех элементов множества $U_n$ сравнимо с произведением всех элементов после умножения на $a$:
\[
    \prod_{x\in U_n} x \equiv \prod_{x\in U_n} ax \pmod n.
\]
Правая часть равна
\[
    a^{\varphi(n)}\prod_{x\in U_n}x.
\]
Так как каждый элемент $x\in U_n$ обратим по модулю $n$, произведение всех элементов также обратимо. Поэтому его можно сократить:
\[
    1\equiv a^{\varphi(n)} \pmod n.
\]
Теорема доказана.
\end{proof}

\subsection{Алгоритм генерации ключей RSA}

\begin{enumerate}
    \item Выбрать два больших различных простых числа $p$ и $q$.
    \item Вычислить
    \[
        n=pq.
    \]
    \item Вычислить
    \[
        \varphi(n)=(p-1)(q-1).
    \]
    \item Выбрать открытый показатель $e$ такой, что
    \[
        1<e<\varphi(n),\qquad \gcd(e,\varphi(n))=1.
    \]
    \item Найти закрытый показатель $d$ как обратный к $e$ по модулю $\varphi(n)$:
    \[
        ed\equiv 1 \pmod{\varphi(n)}.
    \]
    \item Открытый ключ:
    \[
        pk=(n,e).
    \]
    \item Закрытый ключ:
    \[
        sk=(n,d).
    \]
\end{enumerate}

На практике часто используется открытый показатель
\[
    e=65537.
\]

\subsection{Шифрование и расшифрование RSA}

Сообщение представляется числом
\[
    m\in \{0,1,\ldots,n-1\}.
\]

Шифрование:
\[
    c \equiv m^e \pmod n.
\]

Расшифрование:
\[
    m \equiv c^d \pmod n.
\]

\subsection{Корректность RSA}

\begin{theorem}[Корректность RSA]
Пусть $n=pq$, где $p$ и $q$ --- различные простые числа. Пусть числа $e$ и $d$ выбраны так, что
\[
    ed\equiv 1 \pmod{\varphi(n)}.
\]
Тогда для любого сообщения $m\in \mathbb{Z}_n$ выполняется
\[
    (m^e)^d \equiv m \pmod n.
\]
\end{theorem}

\begin{proof}
Из условия
\[
    ed\equiv 1 \pmod{\varphi(n)}
\]
следует, что существует целое $t$ такое, что
\[
    ed = 1+t\varphi(n).
\]

Нужно доказать:
\[
    m^{ed}\equiv m \pmod n.
\]

Так как $n=pq$, достаточно доказать сравнение по модулям $p$ и $q$.

Рассмотрим модуль $p$.

Если $p\mid m$, то
\[
    m\equiv 0 \pmod p.
\]
Тогда
\[
    m^{ed}\equiv 0 \equiv m \pmod p.
\]

Если $p\nmid m$, то $\gcd(m,p)=1$. По малой теореме Ферма:
\[
    m^{p-1}\equiv 1 \pmod p.
\]
Так как
\[
    \varphi(n)=(p-1)(q-1),
\]
число $\varphi(n)$ делится на $p-1$. Следовательно,
\[
    m^{\varphi(n)}\equiv 1 \pmod p.
\]
Тогда
\[
    m^{ed}=m^{1+t\varphi(n)}=m\left(m^{\varphi(n)}\right)^t\equiv m\cdot 1^t\equiv m \pmod p.
\]

Аналогично доказывается:
\[
    m^{ed}\equiv m \pmod q.
\]

Так как $p$ и $q$ взаимно просты, из сравнений
\[
    m^{ed}\equiv m \pmod p,\qquad m^{ed}\equiv m \pmod q
\]
следует
\[
    m^{ed}\equiv m \pmod {pq}.
\]
То есть
\[
    (m^e)^d\equiv m \pmod n.
\]
\end{proof}

\subsection{Минимальный пример RSA}

Возьмём малые числа, непригодные для реальной криптографии:

\[
    p=5,\qquad q=11.
\]
Тогда:
\[
    n=pq=55,
\]
\[
    \varphi(n)=(5-1)(11-1)=4\cdot 10=40.
\]

Выберем
\[
    e=3,
\]
так как
\[
    \gcd(3,40)=1.
\]

Найдём $d$:
\[
    3d\equiv 1 \pmod{40}.
\]
Подходит
\[
    d=27,
\]
поскольку
\[
    3\cdot 27=81\equiv 1 \pmod{40}.
\]

Открытый ключ:
\[
    pk=(55,3).
\]

Закрытый ключ:
\[
    sk=(55,27).
\]

Пусть сообщение
\[
    m=7.
\]

Шифрование:
\[
    c\equiv 7^3 \pmod{55}.
\]
\[
    7^3=343,\qquad 343\bmod 55=13.
\]
Значит,
\[
    c=13.
\]

Расшифрование:
\[
    m\equiv 13^{27}\pmod{55}.
\]
Вычисляя быстрое возведение в степень:
\[
    13^2\equiv 169\equiv 4 \pmod{55},
\]
\[
    13^4\equiv 4^2=16 \pmod{55},
\]
\[
    13^8\equiv 16^2=256\equiv 36 \pmod{55},
\]
\[
    13^{16}\equiv 36^2=1296\equiv 31 \pmod{55}.
\]
Так как
\[
    27=16+8+2+1,
\]
получаем:
\[
    13^{27}\equiv 13^{16}13^8 13^2 13
    \equiv 31\cdot 36\cdot 4\cdot 13 \pmod{55}.
\]
\[
    31\cdot 36=1116\equiv 16 \pmod{55},
\]
\[
    16\cdot 4=64\equiv 9 \pmod{55},
\]
\[
    9\cdot 13=117\equiv 7 \pmod{55}.
\]
Итак,
\[
    13^{27}\equiv 7 \pmod{55}.
\]

\subsection{Практические замечания о RSA}

``Голый'' RSA без дополнения небезопасен:

\begin{enumerate}
    \item Он детерминирован: одинаковые сообщения дают одинаковые шифртексты.
    \item Малые сообщения могут быть уязвимы.
    \item Возможны атаки с выбранным шифртекстом.
\end{enumerate}

Поэтому на практике применяют схемы дополнения:

\begin{itemize}
    \item RSA-OAEP для шифрования;
    \item RSA-PSS для цифровой подписи.
\end{itemize}

\section{Цифровая подпись}

\subsection{Назначение цифровой подписи}

Цифровая подпись обеспечивает:

\begin{enumerate}
    \item целостность сообщения;
    \item аутентификацию подписанта;
    \item неотказуемость.
\end{enumerate}

\begin{definition}
Схема цифровой подписи задаётся тройкой алгоритмов
\[
    (\KeyGen,\Sign,\Ver),
\]
где:
\begin{itemize}
    \item $\KeyGen$ генерирует пару ключей $(pk,sk)$;
    \item $\Sign_{sk}(m)$ создаёт подпись $\sigma$ для сообщения $m$;
    \item $\Ver_{pk}(m,\sigma)$ проверяет подпись.
\end{itemize}
\end{definition}

Корректность:
\[
    \Ver_{pk}(m,\Sign_{sk}(m))=1.
\]

\subsection{RSA-подпись в учебной форме}

В простейшем учебном варианте RSA-подпись работает так.

Пусть:
\[
    pk=(n,e),\qquad sk=(n,d).
\]

Подписание сообщения $m$:
\[
    \sigma \equiv m^d \pmod n.
\]

Проверка:
\[
    \sigma^e \stackrel{?}{\equiv} m \pmod n.
\]

Корректность:
\[
    \sigma^e\equiv (m^d)^e=m^{de}\equiv m \pmod n.
\]

На практике подписывают не само сообщение, а его хеш с использованием специальных схем дополнения:
\[
    \sigma = \Sign_{sk}(H(m)).
\]

\subsection{Минимальный пример RSA-подписи}

Используем параметры из предыдущего примера:
\[
    n=55,\qquad e=3,\qquad d=27.
\]

Пусть сообщение:
\[
    m=7.
\]

Подпись:
\[
    \sigma \equiv 7^{27}\pmod{55}.
\]

Вычислим:
\[
    7^2=49\equiv -6 \pmod{55},
\]
\[
    7^4\equiv (-6)^2=36 \pmod{55},
\]
\[
    7^8\equiv 36^2=1296\equiv 31 \pmod{55},
\]
\[
    7^{16}\equiv 31^2=961\equiv 26 \pmod{55}.
\]

Так как
\[
    27=16+8+2+1,
\]
получаем:
\[
    7^{27}\equiv 7^{16}7^8 7^2 7
    \equiv 26\cdot 31\cdot 49\cdot 7 \pmod{55}.
\]
\[
    26\cdot 31=806\equiv 36 \pmod{55},
\]
\[
    36\cdot 49=1764\equiv 4 \pmod{55},
\]
\[
    4\cdot 7=28 \pmod{55}.
\]
Значит,
\[
    \sigma=28.
\]

Проверка:
\[
    \sigma^e=28^3\pmod{55}.
\]
\[
    28^2=784\equiv 14 \pmod{55},
\]
\[
    28^3\equiv 14\cdot 28=392\equiv 7 \pmod{55}.
\]
Проверка успешна:
\[
    \sigma^e\equiv m \pmod{55}.
\]

\section{Методы генерации и распределения ключей}

\subsection{Проблема распределения ключей}

В симметрической криптографии для $N$ участников, каждому из которых нужна отдельная защищённая связь с каждым другим, может потребоваться
\[
    \frac{N(N-1)}{2}
\]
различных ключей.

Это плохо масштабируется.

Асимметрическая криптография уменьшает проблему: каждый пользователь хранит один закрытый ключ и публикует один открытый ключ. Но появляется задача проверки подлинности открытого ключа.

\subsection{Генерация симметрических ключей}

Секретные ключи должны генерироваться с помощью криптографически стойкого генератора случайных чисел.

Основные требования:

\begin{enumerate}
    \item высокая энтропия;
    \item непредсказуемость;
    \item отсутствие повторного использования в неподходящих режимах;
    \item достаточная длина ключа.
\end{enumerate}

Типичные длины современных симметрических ключей:
\[
    128,\ 192,\ 256 \text{ бит}.
\]

\subsection{Генерация RSA-ключей}

Алгоритм генерации RSA-ключей на практике:

\begin{enumerate}
    \item Сгенерировать случайное нечётное число нужной длины.
    \item Проверить его на простоту вероятностными тестами, например Миллера--Рабина.
    \item Повторить до получения простого $p$.
    \item Аналогично получить простое $q$, причём $p\ne q$.
    \item Проверить дополнительные условия безопасности:
    \begin{itemize}
        \item $p$ и $q$ не должны быть слишком близки;
        \item $p-1$ и $q-1$ не должны иметь слишком специальную структуру;
        \item $\gcd(e,\varphi(n))=1$.
    \end{itemize}
    \item Вычислить $n=pq$ и закрытый показатель $d$.
\end{enumerate}

\subsection{Тест Миллера--Рабина}

Тест Миллера--Рабина --- вероятностный тест простоты.

Пусть проверяется нечётное число $n>2$.

Представим:
\[
    n-1=2^s d,
\]
где $d$ нечётно.

Для случайного основания $a$, где
\[
    2\le a\le n-2,
\]
вычисляется:
\[
    x=a^d\bmod n.
\]

Если
\[
    x=1
\]
или
\[
    x=n-1,
\]
то раунд теста пройден.

Иначе последовательно вычисляют:
\[
    x\leftarrow x^2\bmod n
\]
до $s-1$ раз. Если на каком-то шаге
\[
    x=n-1,
\]
раунд пройден. Если нет, число составное.

\subsubsection*{Алгоритм Миллера--Рабина}

\begin{enumerate}
    \item Если $n=2$ или $n=3$, вернуть ``простое''.
    \item Если $n<2$ или $n$ чётно, вернуть ``составное''.
    \item Представить $n-1=2^s d$, где $d$ нечётно.
    \item Повторить $t$ раз:
    \begin{enumerate}
        \item выбрать случайное $a\in\{2,\ldots,n-2\}$;
        \item вычислить $x=a^d\bmod n$;
        \item если $x=1$ или $x=n-1$, перейти к следующему раунду;
        \item повторить $s-1$ раз:
        \[
            x\leftarrow x^2\bmod n;
        \]
        если $x=n-1$, перейти к следующему раунду;
        \item вернуть ``составное''.
    \end{enumerate}
    \item Вернуть ``вероятно простое''.
\end{enumerate}

\subsection{Протокол Диффи--Хеллмана}

Протокол Диффи--Хеллмана позволяет двум сторонам выработать общий секрет по открытому каналу.

\subsubsection*{Параметры}

Пусть задана циклическая группа $G$ порядка $q$ с образующим $g$.

\subsubsection*{Алгоритм}

\begin{enumerate}
    \item Алиса выбирает случайное секретное число
    \[
        a\in \{1,\ldots,q-1\}
    \]
    и отправляет Бобу
    \[
        A=g^a.
    \]
    \item Боб выбирает случайное секретное число
    \[
        b\in \{1,\ldots,q-1\}
    \]
    и отправляет Алисе
    \[
        B=g^b.
    \]
    \item Алиса вычисляет
    \[
        K=B^a=(g^b)^a=g^{ab}.
    \]
    \item Боб вычисляет
    \[
        K=A^b=(g^a)^b=g^{ab}.
    \]
\end{enumerate}

В итоге обе стороны получают один и тот же секрет:
\[
    K=g^{ab}.
\]

\subsubsection*{Минимальный пример}

Возьмём малые учебные параметры:
\[
    p=23,\qquad g=5.
\]

Алиса выбирает:
\[
    a=6.
\]
Она вычисляет:
\[
    A=5^6\bmod 23.
\]
\[
    5^2=25\equiv 2,\qquad 5^4\equiv 4,\qquad 5^6\equiv 4\cdot 2=8.
\]
Значит,
\[
    A=8.
\]

Боб выбирает:
\[
    b=15.
\]
Он вычисляет:
\[
    B=5^{15}\bmod 23=19.
\]

Алиса вычисляет:
\[
    K=B^a=19^6\bmod 23.
\]
\[
    19^2=361\equiv 16,\qquad 19^4\equiv 16^2=256\equiv 3,
\]
\[
    19^6\equiv 19^4\cdot 19^2\equiv 3\cdot 16=48\equiv 2.
\]

Боб вычисляет:
\[
    K=A^b=8^{15}\bmod 23.
\]
\[
    8^2=64\equiv 18,\quad
    8^4\equiv 18^2=324\equiv 2,
\]
\[
    8^8\equiv 2^2=4.
\]
Так как
\[
    15=8+4+2+1,
\]
\[
    8^{15}\equiv 8^8\cdot 8^4\cdot 8^2\cdot 8
    \equiv 4\cdot 2\cdot 18\cdot 8.
\]
\[
    4\cdot 2=8,\quad 8\cdot 18=144\equiv 6,\quad 6\cdot 8=48\equiv 2.
\]
Обе стороны получили:
\[
    K=2.
\]

\subsection{Уязвимость Диффи--Хеллмана к атаке посредника}

Обычный Диффи--Хеллман без аутентификации уязвим к атаке ``человек посередине''.

Злоумышленник может подменить $A$ и $B$, установив один ключ с Алисой и другой ключ с Бобом. Поэтому в реальных протоколах Диффи--Хеллман должен аутентифицироваться, например цифровыми подписями или сертификатами.

\subsection{Инфраструктура открытых ключей}

Для проверки принадлежности открытого ключа конкретному субъекту используется PKI --- инфраструктура открытых ключей.

Основные элементы PKI:

\begin{enumerate}
    \item \textbf{Удостоверяющий центр} --- организация, выпускающая сертификаты.
    \item \textbf{Сертификат открытого ключа} --- подписанный документ, связывающий открытый ключ с владельцем.
    \item \textbf{Цепочка сертификации} --- последовательность сертификатов от конечного пользователя к доверенному корневому центру.
    \item \textbf{Списки отзыва} или онлайн-проверка статуса сертификата --- механизмы проверки, не был ли сертификат отозван.
\end{enumerate}

\section{Типичные атаки и практические замечания}

\subsection{Атаки на шифры}

Основные типы атак:

\begin{enumerate}
    \item \textbf{Атака только по шифртексту}: известны только шифртексты.
    \item \textbf{Атака с известным открытым текстом}: известны пары $(m,c)$.
    \item \textbf{Атака с выбранным открытым текстом}: противник может получать шифртексты для выбранных сообщений.
    \item \textbf{Атака с выбранным шифртекстом}: противник может получать расшифрования выбранных шифртекстов.
    \item \textbf{Побочные каналы}: время выполнения, энергопотребление, электромагнитные излучения, ошибки вычислений.
\end{enumerate}

\subsection{Режим ECB}

Режим ECB шифрует каждый блок независимо:
\[
    c_i=E_k(m_i).
\]

Недостаток: одинаковые открытые блоки дают одинаковые шифртексты. Поэтому структура сообщения частично сохраняется.

ECB не рекомендуется для шифрования длинных сообщений.

\subsection{Режим CBC}

В режиме CBC:
\[
    c_0=IV,
\]
\[
    c_i=E_k(m_i\oplus c_{i-1}).
\]

Для расшифрования:
\[
    m_i=D_k(c_i)\oplus c_{i-1}.
\]

Требуется случайный и непредсказуемый вектор инициализации $IV$.

\subsection{Режим CTR}

В режиме CTR блочный шифр используется как генератор гаммы:
\[
    z_i=E_k(\text{nonce}\,\|\,i),
\]
\[
    c_i=m_i\oplus z_i.
\]

Важно: нельзя повторно использовать одну и ту же пару $(k,\text{nonce})$.

\subsection{Аутентифицированное шифрование}

На практике желательно использовать не просто шифрование, а аутентифицированное шифрование, которое одновременно обеспечивает:

\begin{enumerate}
    \item конфиденциальность;
    \item целостность;
    \item аутентичность.
\end{enumerate}

Примеры:
\begin{itemize}
    \item AES-GCM;
    \item ChaCha20-Poly1305;
    \item режимы на основе ГОСТ-алгоритмов с имитовставкой.
\end{itemize}

\section{Сравнительная таблица основных механизмов}

\begin{center}
\begin{tabular}{llll}
\toprule
Механизм & Ключи & Основная цель & Пример \\
\midrule
Симметрическое шифрование & общий секретный & конфиденциальность & AES, Магма \\
Асимметрическое шифрование & открытый/закрытый & конфиденциальность & RSA-OAEP \\
MAC & общий секретный & целостность, аутентичность & HMAC \\
Цифровая подпись & открытый/закрытый & целостность, авторство & RSA-PSS, ГОСТ Р 34.10 \\
Хеш-функция & ключ не нужен & отпечаток данных & SHA-256, Стрибог \\
Обмен ключами & временные секреты & общий ключ & Diffie--Hellman \\
\bottomrule
\end{tabular}
\end{center}

\section{Краткие выводы}

\begin{enumerate}
    \item Конфиденциальность достигается шифрованием, целостность --- MAC или цифровой подписью.
    \item Совершенная секретность возможна, но требует ключа длиной не меньше сообщения.
    \item Современная криптография в основном использует вычислительную стойкость.
    \item DES исторически важен, но практически устарел из-за 56-битного ключа.
    \item AES является основным международным стандартом блочного шифрования.
    \item Российские современные блочные шифры --- Магма и Кузнечик.
    \item RSA основан на трудности факторизации больших чисел.
    \item Цифровая подпись обеспечивает целостность, аутентификацию и неотказуемость.
    \item Распределение ключей решается с помощью протоколов обмена ключами и инфраструктуры открытых ключей.
\end{enumerate}

\end{document}
